\section{The Instant and Eternity}

\begin{quotex}
This is an article by \textbf{Guido de Giorgio}, titled \textit{The Instant and Eternity}, first published in Diorama Filosofico in 1939.
\end{quotex}

\paragraph{The Myth of the Future}
We can say that the sacred is distinguished from the profane in what is essentially oriented toward the past to fix the stages of a procession which necessarily finds its culmination in a ``present''. This ``present'' is the metaphysical point where eternity throws itself, where the worlds dissolve in a fullness without limits, a duration without rhythm, a bliss without end. The present is eternity, the past is only the vestibule that leads toward, which inserts into eternity. To repeat, to retrace the whole cycle which is realized in the point means to carry with oneself the experience of the centuries, all cosmic evolution to unravel the framework in the pupil of God.

Faust could not stop the instant because he knew only the caducity of the instance, the immediate iridescence of illusion, the vertigo which submerges instead of transfiguring, the ``shadow of the flesh'', the labile and evanescent phantasm, not that which in God resides an infinite momentariness which is the mystery of the eternal now. Such are the two aspects of the ``instant'', according to which one places himself on the human or divine plane; it is about two apparently opposed and divergent points which mark two worlds, two rhythms, two realities, of which one is absolute, true, the other fallacious and illusory. Faust's words ``stop, you are so beautiful!'' is only a not very original lyrical substitute in the face of the unfathomed plenitude of the Ineffable where the mystery of the divine gestation takes place. The myth of purification through aesthetics is only the very fragile bridge thrown by modern imbecility onto the momentariness of the human-cosmic illusion in order to evade the positive certitude of the mystery, an impassable wall or else by the dizzying passing of the wing, i.e., of the Spirit of God.

That is why the modern world oscillates between a dead past and a nebulous future, between what is no longer and what will never be, except in the hope that anticipates and builds. Traditional wisdom, on the contrary, turns toward the past, lives it, enriches it, \emph{updates} it, inserting itself in it to lead it fully into the present and renew it in the \emph{ver aeternum} eternal spring which the Ancients attributed to the Golden Age, pointing to the perennial germination of the Truth, the multitude of transfiguring states, life which knows neither birth nor death, for it uncoils itself in the bliss of the realizing consciousness. But for moderns the past is past, dead, finished, concluded, closed, irremediable: ``le déjà vu, le déjà vécu'' already seen, already lived , says Bergson, according to a psychological orientation that clearly manifests all the nostalgic sentimentality of the small man horribly enslaved by his small world. So that between a dead past and a future not yet born, the twilight present swings, at once a cloudy waning and a too pale dawn, in sum a veritable pause in agony. And from this erroneous vision the \emph{myth of the future} is derived, the tension toward what is not, toward what will never be because in reality only the present, in absorbing the past, is the dynamic point, \emph{the whole bow of the ship} which faces the horizon \emph{but never reaches it}.


\paragraph{The Past as Death and Life}

\begin{quotationx}
This is the second section of \textbf{Guido de Giorgio}`s essay The Instant and Eternity. Here de Giorgio compares modern and traditional man. He brings in many themes previously covered here: the continuity of the past; the notion of depth in the flow of history, i.e., ``subterranean history''; the past understood in the interiority, not as something external. Moderns can never rest, they can never ``be'', because they are always becoming, looking to the future to pass judgment on them. But the moderns of the future will despise them as much as the moderns of today despise their own roots, heritage and past.

Burying beetles bury the carcasses of small vertebrates such as birds and rodents as a food source for their larvae. This is the perfect description of modern man as a beetle. Although the essay was written in 1939, it is even more applicable today.
\end{quotationx}


Modern man can be compared to a burying beetle nécrophore that longs for the day that never comes: the cadaver that he carries is the past, his inert, sterile heritage, and the day that he awaits is the future, the imaginary descent, the glorious completion of a chimeric unfinished childbirth. We will remark that all modern men, the ``great men of history'', wait for a definitive judgment of the future on their oeuvre, for perhaps they feel, consciously or not, that nothing of what they made relates traditionally to the royal stream of the past, not is it capable of resisting the movements of the magnetic needle of the present, the brief and momentary instant having an impact on many abysses other than the insignificant trace of the passing cloud. That is why ancient man is a bearer of worlds: he did not leave the past behind him, but harvests it and carries it along, in such a way to construct in reality a single incidental point, the sole present, the current time, while modern man, discarding a heavy burden for his not very virile shoulders, is light, inconsistent, and through fear of being thrown to the earth by the blows of the crosswinds, clings to the machine, which is both his cradle and his tomb. For the myth of the future is associated with the myth of speed which---if we consider its function, its interior plan---is the abolition of the past in the already traversed, the imperceptibility of the present minimized in the permanent expectation of the future. The readers who would like to deepen these insights in a penetrating manner will find more than one way leading easily to the comprehension of some major truths. We desire to establish here, with a certain insistence, only some critical reflections whose development perspective will turn out to be clearer and surer.

We therefore understand that modern man and ancient man are absolutely opposed and like the antipodes, in the literal sense of the term, the one by relationship to the other: tied to a same line of descent but turned toward different heavens with different constellations, although the same impassible sun throws light on that line of descent in what for one is day for the other, night. For the Ancients, in effect, the past is everything, for the moderns, nothing, even when they have the illusion of absent mindedly looking in the past for solutions to the questions of the present---what are called the ``warnings'', the ``teachings'' of the past--, around sentimental fantasies exploited with a cynical opportunism according to the circumstances and proposed to the credulity of the naïve for the most pitiable perpetrations. Rhetoric, which triumphs today as never before in today's cloudy and swampy Europe, has recourse to the most bestial ruses to obtain the consent of the plebes as the audience and makes use of the past as a remedy against all illnesses, the universal balm, the supports of the present, but of a momentary usage as if to ward off the \emph{Vae soli}! Woe to him who is alone, Eccl. 4:10 Modern man, in reality, is already frozen in the past, no longer lives it, and takes from it only dust and ruin: he studies it, classifies it, ignores it. The more detailed the inquiry is made, the more it becomes sketchy, each one seeking subsequently to make life blow over these bones asleep in the slumber of death. Thus, when they turn toward the past to study it, the moderns then succumb to the same illusion as when they believe, for example, that the photograph is closer than the truth, although it is denatured totally in being fixing in the momentariness of something already passed. But independently of study, let us see if the moderns use the past according to life. Who says past, says tradition, i.e., interior, dynamic unification, not exterior adhesion, not opportunistic sympathy, not simple position or situation; in other words, there should be between the past and the present, \emph{continuity}, \emph{immutability}, or, better said, a rhythmic development so regular, permanent, internal, that it would appear to be indifferent. Antiquity, in fact, is characterized by a constant tonality which endures, immobile, from one epoch to the other; there is and must be a change, but it takes place in depth, in the interior strata, invisibly, we are tempted to say, in a way to not disrupt the regularity of the rhythm.


\paragraph{The Meaning of the Past}

\begin{quotex}
This is third section of an article by \textbf{Guido de Giorgio}, titled The Instant and Eternity, first published in Diorama Filosofico in 1939.

I am constantly struck by so many people today who only perceive darkness, ignorance, or intolerance in the past. They tell me about dictators oppressively their people, who must be saved; some unfortunates who cannot afford this or that without a new government program; some marginal group that must be accepted; and so on. They long for some mythical future when all will be well with the world. Yet they do nothing for their own family, their own children and grandchildren; they assume their current life will continue, just as it is now. They cannot conceive that someday they will be oppressed, impoverished, or marginalized, with nothing to hope for. Those who cut their ties to the past don't have a future, because they don't even know what it is.
\end{quotex}


We said many times that ancient cultures are immobile or seem to be such; but that is precisely their greatness, that fundamental stability that removes all contrasts, integrates all rhythms into the central vein, to the \emph{traditional type}, who alone remains in the wholeness of its formative efficacy. That is why whoever intends to remain in the pure domain of traditional truth, always turns, logically, toward the past, to retrace the stages of certitude and add them to his experience. This, under that angle, is then recapitulative and conclusive: it is not an exterior repetition, but grants its rhythm to that which is none other than is own face, today still ignored, from now on found again and vivified. It is very difficult to explain certain things to those who hold onto dualist positions and who think that there is something besides the Truth, which is God eternally present. Truth: there only can \emph{one become what he is}, i.e., one transcends the sphere of human limitations to live the same beat of the infinite.

When we say \emph{ancient} we mean everything that is valuable, perennial, traditionally authentic in the past of the East and the West, whether it is about a distant or near, doctrinal or poetic, past: it matters little, provided that it reflects, in the truth of its expression, the great light of the Superworld. Beyond the Sacred Books, there is Poetry and sacred Art. There is finally all the forms of activity which, in the past, always relate to a truth of a higher order, be it in a modest utensil, and in the fabrication and the destination of objects of current usage. That past, as we mean it, and all those who are searching only for the truth of God should mean it, is life, creative rhythm, an inexhaustible deposit of wisdom which renews itself each time it is actualized by a new experience. But it is especially the reality of a vibrant life because vivified by the perennial breath of traditional energy. Moderns consider the past as a relic whose ruins they borrow and around which they prowl with the curiosity of photographers and archeologists: who, among them, accepts the past totally, assumes it in all its fullness, not in order to seize form it some fragments and exalt them, but in order to integrate it in one's experience of life while recapitulating it in a creative manner?


\paragraph{The Tangible Immediacy of the Present}

\begin{quotex}
This is final section of an article by \textbf{Guido de Giorgio}, titled The Instant and Eternity, first published in Diorama Filosofico in 1939.
\end{quotex}


How many of Dante's admirers are there who are not content in glorifying his verses or expression---something absolutely exterior and superficial---but who in applying the doctrine, the knowledge on all the planes of being to which they relate and in the totality of the Heavenly Voyage.

The past is nothing if it is not integrated, lived, validated by personal experience, by life, if it is not totalized and exalted in the great shaking of the eternal now. The moderns, when they are not fornicating in the past like the thieves in a necropolis, turn their backs on it, contemplating then the hypothetical ``sun of the future'' which will never shine, because the future exists only as the last evanescent border of an arduous vision, a mirage and nothing more, a fallacious projection stained by the spasms of their own insufficiency. The ``non-achievement'' in the face of the Truth, the incurable sentiment of one who does not know nor wants to know, does not know how to nor wants to carry with himself all the weight of the world, to assume it in the divine instant, created the myth of the future. Turning obstinately their backs to that which is, they wait with curiosity for what is not, for what will be, and they long for the confirmation of a dream by an illusory reflection of the dream itself, in a nocturnal march of phantoms that alone engender the present, by the spontaneity of its flux and its mirage. Strange speculation on the future, which makes them forget the treasures of the past and the tangible immediacy of the present. For they are really only in the present, with all worlds, in the essential unity of the point, jewel of all jewels, the eternal eye of God!

We would like to still say some other things, but we prefer to conclude with the words of \textbf{Zarathustra}:

\begin{quotex}
To these men of today will I not be light, nor be called light. Them, will I blind with the lightning of my wisdom! Put out their eyes!
\flright{\textsc{Friedrich Nietzsche}, \emph{Zarathustra}}
\end{quotex}



\flrightit{Posted on 2012-06-12 by Cologero}

\begin{center}* * *\end{center}

\begin{footnotesize}
\begin{sffamily}
\texttt{Egemen on 2012-06-12 at 04:51 said:}

A nice dealing about the sacred past of civilization. I want to share a song and poem of one of the esoteric orders of the that is bekta?ilik (\url{http://en.wikipedia.org/wiki/Bektashi\_Order}), (\url{http://it.wikipedia.org/wiki/Bektashi})


And this is the youtube link of song;
\url{http://www.youtube.com/watch?v=m-aWLCGLCNo}

And these are the lyrics which will translate and explain by me later... 
\begin{verse}
Ondort bin yil gezdim pervanelikte\\
Sidk-i ismin buldum divanelikte\\
?ctim sarabini mestanelikte\\
Kirklarin ceminde dara düs oldum\\
Guruh-u naci'ye özümü kattim\\
?nsan s?fatindan cok geldim gittim\\
Bülbül oldum firdevs baginda öttüm\\
Bir zaman gül icin zara düs oldum
\end{verse}

I had travelled in space for fourteen thousand years (pervanelik means space or outer space and also pervanelik means planets; if a object turning around another thing, you can say for object; it is the pervane of other thing)

I had found the true meaning of your name (sidk-i ism) in madness (divanelik means madness or extacy)

I had drunk your wine at mestanelik (mestanelik means a place which people is getting transcendence state of mind)

I had forced to be brought to book at forty's court (in ancient asian tribes or little orders there was a court that call sometimes ``gathering'', which judges some problems of society. Forty's are prophets from Adam to Muhammed)

I had put in the forty's to my essence (güruh-u naci is the persian idiom of Forty's)

I had come and gone from the presence of human (or Adam Kadmon), (that means i had many experience of transfiguraiton, transformation or metamorphism... )

I had had form of nightingale and sing at the vineyard of Firdevs... (Firdevs is the highest level of the heaven; maybe as mentioned in Dante's Comedia)

i had fall to a garden for rose, for a moment... After all this explanations, there is a little story about origins of that song:

Before Islam and other monotheistic religions, ancient central asian nations had their pantheons, their celestial beings as elohim in torah. One of them called as ``Gabriel, Cebrail or Cibril, most probably after influence of Islam, and this Gabriel is humanlike creature of the god.

This human's breed was in danger, because of their green sun which Gabriel's home planet rounds around it (pervanelik), had come to hers ver end of life.

If Gabriel could not hade found a way to rescue his breed, they will have been destroyed with their planet. Hereby, Gabriel or Cibril had get a quest for a new planet to live on. And Cibril had prepared his spacecraft and hit on the roads of outer space, after that he had met with the god. And god asked to the Cibril:

``Who are you and who am I?''

Cibril answered was a little bit defiant;

``You are You and I am me''

God had not liked that answer and He punished Cibril by rampant traveling in outer space about fourteen thousand years.

In this punisment era, Cibril was getting into the madness and suddenly God appears and ask his same question again:

``Who are you and who am I?''

After all these hard years, Cibril had come to his senses and answers:

``You are the Creator, I am the creation''

After being excused, Cibril encounters a blue planet, they land to that planet and search of the lands and waters. They recognize, a creature that has similar genetically similar to them is living that blue planet. They decide to make a mergence between their breed and that creature. They had get some of these creatures with them and return their home planet.

At the Forty's Court; they had made some genetical experiments about mergence of breeds and they found a new species that are similar to them. After that, the creators of human generation had get that new species and return to the blue planet and so on as you know.

P.S. I listened that story in 1987 from a Bektasi person. After several years, I met with Zecharia Sitchin and his Works

\url{http://www.sitchin.com/}


After some research about that story, I found it is not only belongs to Bektashi order, and you can find similar stories in Dharma and Ancient Chinese Traditions... 

\hfill

\end{sffamily}
\end{footnotesize}


